\section{Biological Example Family IV: Bistable and Sequential Regimes}

Many neural and behavioural systems exhibit either bistability or sequential
regime structure. Bistable systems matter because the same system may support
two relatively stable outcomes separated by a threshold. Sequential regimes
matter because order can become part of the system architecture.

The self-organizing attractor paper already indexed for this project is
valuable here. Its medium-strength contribution is not merely that attractors
exist, but that repeated sequence exposure can reshape coupling structure and
stabilize later sequential behaviour.

\begin{example}[Study sequence as learned transition graph]
Suppose a reader repeatedly performs the same sequence:
\begin{enumerate}[nosep]
  \item sit at the same desk,
  \item open the same notebook,
  \item review the last unresolved line,
  \item begin with one worked problem.
\end{enumerate}
Over time, the sequence itself can become easier to enter and traverse. The
mathematical value of sequential-dynamics language is that the target is no
longer only a state; it can also be a \emph{path} through state space.
\end{example}
